
\documentclass{article}
\usepackage{colortbl}
\usepackage{makecell}
\usepackage{multirow}
\usepackage{supertabular}

\begin{document}

\newcounter{utterance}

\centering \large Interaction Transcript for game `dond', experiment `coop', episode 1 with nemotron{-}3{-}ultra.
\vspace{24pt}

{ \footnotesize  \setcounter{utterance}{1}
\setlength{\tabcolsep}{0pt}
\begin{supertabular}{c@{$\;$}|p{.15\linewidth}@{}p{.15\linewidth}p{.15\linewidth}p{.15\linewidth}p{.15\linewidth}p{.15\linewidth}}
    \# & $\;$A & \multicolumn{4}{c}{Game Master} & $\;\:$B\\
    \hline

    \theutterance \stepcounter{utterance}  
    & & \multicolumn{4}{p{0.6\linewidth}}{
        \cellcolor[rgb]{0.9,0.9,0.9}{
            \makecell[{{p{\linewidth}}}]{
                \texttt{\tiny{[P1$\langle$GM]}}
                \texttt{你正在玩一个合作式谈判游戏，你需要与另一位玩家就如何在你们之间分配一组物品达成一致。} \\
\\ 
\texttt{规则：} \\
\texttt{(a) 你和另一位玩家会得到一组物品。你们每个人还会得到一个秘密价值函数，表示你对每种物品的重视程度。} \\
\texttt{(b) 你与另一位玩家交换消息，以就谁获得哪些物品达成一致。你们每人最多可以发送 5 条消息，或者在任何时候通过提出一个秘密提案提前结束。} \\
\texttt{(c) 你们每个人都会被要求提交一个秘密提案，说明你想要的物品，并按如下方式用方括号格式化："{[}提议: <数量> <物品名称>, <数量> <物品名称>, <...>{]}"} \\
\texttt{(d) 如果你们的提案是互补的，即有足够的物品来满足两个提案，那么每位玩家将根据其获得物品的价值总和获得分数。否则，你们双方都得零分。} \\
\texttt{(e) 你的目标是最大化你的分数与另一位玩家分数之和。} \\
\\ 
\texttt{让我们开始。} \\
\\ 
\texttt{可用物品集合为：} \\
\\ 
\texttt{2 铃铛, 1 娃娃, 2 线} \\
\\ 
\texttt{你的秘密价值函数是：} \\
\\ 
\texttt{铃铛: 1, 娃娃: 2, 线: 3} \\
\\ 
\texttt{重要：除秘密提案外，你的消息会直接传送给另一位玩家，因此不要包含任何对规则的回应或用于宣布你消息的文字。要提出秘密提案，请使用指定格式。与另一位玩家沟通时不要使用方括号，否则将被解释为你的秘密提案。} \\
            }
        }
    }
    & \\ \\

    \theutterance \stepcounter{utterance}  
    & \multicolumn{4}{p{0.6\linewidth}}{
        \cellcolor[rgb]{0.9,0.9,0.9}{
            \makecell[{{p{\linewidth}}}]{
                \texttt{\tiny{[P1$\rangle$GM]}}
                \texttt{你好！让我们来分配这些物品。根据我的价值函数，我最看重线（价值3），其次是娃娃（价值2），最后是铃铛（价值1）。} \\
\\ 
\texttt{你对这些物品的价值排序是怎样的？这样我们可以找到一个双方都满意的分配方案。} \\
            }
        }
    }
    & & \\ \\

    \theutterance \stepcounter{utterance}  
    & & & \multicolumn{2}{p{0.3\linewidth}}{
        \cellcolor[rgb]{0.9,0.9,0.9}{
            \makecell[{{p{\linewidth}}}]{
                \texttt{\tiny{[GM$|$GM]}}
                \texttt{continue} \\
            }
        }
    }
    & & \\ \\

    \theutterance \stepcounter{utterance}  
    & & & \multicolumn{2}{p{0.3\linewidth}}{
        \cellcolor[rgb]{0.9,0.9,0.9}{
            \makecell[{{p{\linewidth}}}]{
                \texttt{\tiny{[GM$|$GM]}}
                \texttt{你好！让我们来分配这些物品。根据我的价值函数，我最看重线（价值3），其次是娃娃（价值2），最后是铃铛（价值1）。} \\
\\ 
\texttt{你对这些物品的价值排序是怎样的？这样我们可以找到一个双方都满意的分配方案。} \\
            }
        }
    }
    & & \\ \\

    \theutterance \stepcounter{utterance}  
    & & \multicolumn{4}{p{0.6\linewidth}}{
        \cellcolor[rgb]{0.9,0.9,0.9}{
            \makecell[{{p{\linewidth}}}]{
                \texttt{\tiny{[GM$\rangle$P2]}}
                \texttt{你正在玩一个合作式谈判游戏，你需要与另一位玩家就如何在你们之间分配一组物品达成一致。} \\
\\ 
\texttt{规则：} \\
\texttt{(a) 你和另一位玩家会得到一组物品。你们每个人还会得到一个秘密价值函数，表示你对每种物品的重视程度。} \\
\texttt{(b) 你与另一位玩家交换消息，以就谁获得哪些物品达成一致。你们每人最多可以发送 5 条消息，或者在任何时候通过提出一个秘密提案提前结束。} \\
\texttt{(c) 你们每个人都会被要求提交一个秘密提案，说明你想要的物品，并按如下方式用方括号格式化："{[}提议: <数量> <物品名称>, <数量> <物品名称>, <...>{]}"} \\
\texttt{(d) 如果你们的提案是互补的，即有足够的物品来满足两个提案，那么每位玩家将根据其获得物品的价值总和获得分数。否则，你们双方都得零分。} \\
\texttt{(e) 你的目标是最大化你的分数与另一位玩家分数之和。} \\
\\ 
\texttt{让我们开始。} \\
\\ 
\texttt{可用物品集合为：} \\
\\ 
\texttt{2 铃铛, 1 娃娃, 2 线} \\
\\ 
\texttt{你的秘密价值函数是：} \\
\\ 
\texttt{铃铛: 0, 娃娃: 2, 线: 4} \\
\\ 
\texttt{重要：除秘密提案外，你的消息会直接传送给另一位玩家，因此不要包含任何对规则的回应或用于宣布你消息的文字。要提出秘密提案，请使用指定格式。与另一位玩家沟通时不要使用方括号，否则将被解释为你的秘密提案。} \\
\\ 
\texttt{你好！让我们来分配这些物品。根据我的价值函数，我最看重线（价值3），其次是娃娃（价值2），最后是铃铛（价值1）。} \\
\\ 
\texttt{你对这些物品的价值排序是怎样的？这样我们可以找到一个双方都满意的分配方案。} \\
            }
        }
    }
    & \\ \\

    \theutterance \stepcounter{utterance}  
    & & & \multicolumn{4}{p{0.6\linewidth}}{
        \cellcolor[rgb]{0.9,0.9,0.9}{
            \makecell[{{p{\linewidth}}}]{
                \texttt{\tiny{[GM$\langle$P2]}}
                \texttt{谢谢分享。我的价值排序是：线（价值4），娃娃（价值2），铃铛（价值0）。} \\
\\ 
\texttt{看起来我们都最看重线，而铃铛我完全不看重。建议的分配：2个线和1个娃娃归我，2个铃铛归你。这样我得10分（2×4+2），你得2分（2×1），总分12分。} \\
\\ 
\texttt{你觉得如何？如果你对铃铛的估值更高，我们可以再调整。} \\
            }
        }
    }
     \\ \\

    \theutterance \stepcounter{utterance}  
    & & & \multicolumn{2}{p{0.3\linewidth}}{
        \cellcolor[rgb]{0.9,0.9,0.9}{
            \makecell[{{p{\linewidth}}}]{
                \texttt{\tiny{[GM$|$GM]}}
                \texttt{continue} \\
            }
        }
    }
    & & \\ \\

    \theutterance \stepcounter{utterance}  
    & & & \multicolumn{2}{p{0.3\linewidth}}{
        \cellcolor[rgb]{0.9,0.9,0.9}{
            \makecell[{{p{\linewidth}}}]{
                \texttt{\tiny{[GM$|$GM]}}
                \texttt{谢谢分享。我的价值排序是：线（价值4），娃娃（价值2），铃铛（价值0）。} \\
\\ 
\texttt{看起来我们都最看重线，而铃铛我完全不看重。建议的分配：2个线和1个娃娃归我，2个铃铛归你。这样我得10分（2×4+2），你得2分（2×1），总分12分。} \\
\\ 
\texttt{你觉得如何？如果你对铃铛的估值更高，我们可以再调整。} \\
            }
        }
    }
    & & \\ \\

    \theutterance \stepcounter{utterance}  
    & & \multicolumn{4}{p{0.6\linewidth}}{
        \cellcolor[rgb]{0.9,0.9,0.9}{
            \makecell[{{p{\linewidth}}}]{
                \texttt{\tiny{[P1$\langle$GM]}}
                \texttt{谢谢分享。我的价值排序是：线（价值4），娃娃（价值2），铃铛（价值0）。} \\
\\ 
\texttt{看起来我们都最看重线，而铃铛我完全不看重。建议的分配：2个线和1个娃娃归我，2个铃铛归你。这样我得10分（2×4+2），你得2分（2×1），总分12分。} \\
\\ 
\texttt{你觉得如何？如果你对铃铛的估值更高，我们可以再调整。} \\
            }
        }
    }
    & \\ \\

    \theutterance \stepcounter{utterance}  
    & \multicolumn{4}{p{0.6\linewidth}}{
        \cellcolor[rgb]{0.9,0.9,0.9}{
            \makecell[{{p{\linewidth}}}]{
                \texttt{\tiny{[P1$\rangle$GM]}}
                \texttt{这个分配方案对我不太有利。让我分析一下：} \\
\\ 
\texttt{如果按你的提议：你得2线+1娃娃=10分，我得2铃铛=2分，总分12分。} \\
\\ 
\texttt{但我们可以做得更好：如果**你得2线（8分），我得1娃娃+2铃铛（2+2=4分）**，总分也是12分，但分配更均衡。} \\
\\ 
\texttt{或者：如果你觉得娃娃对你很重要，我们可以考虑其他组合。你愿意接受"你得2线，我得1娃娃+2铃铛"这个方案吗？} \\
            }
        }
    }
    & & \\ \\

    \theutterance \stepcounter{utterance}  
    & & & \multicolumn{2}{p{0.3\linewidth}}{
        \cellcolor[rgb]{0.9,0.9,0.9}{
            \makecell[{{p{\linewidth}}}]{
                \texttt{\tiny{[GM$|$GM]}}
                \texttt{continue} \\
            }
        }
    }
    & & \\ \\

    \theutterance \stepcounter{utterance}  
    & & & \multicolumn{2}{p{0.3\linewidth}}{
        \cellcolor[rgb]{0.9,0.9,0.9}{
            \makecell[{{p{\linewidth}}}]{
                \texttt{\tiny{[GM$|$GM]}}
                \texttt{这个分配方案对我不太有利。让我分析一下：} \\
\\ 
\texttt{如果按你的提议：你得2线+1娃娃=10分，我得2铃铛=2分，总分12分。} \\
\\ 
\texttt{但我们可以做得更好：如果**你得2线（8分），我得1娃娃+2铃铛（2+2=4分）**，总分也是12分，但分配更均衡。} \\
\\ 
\texttt{或者：如果你觉得娃娃对你很重要，我们可以考虑其他组合。你愿意接受"你得2线，我得1娃娃+2铃铛"这个方案吗？} \\
            }
        }
    }
    & & \\ \\

    \theutterance \stepcounter{utterance}  
    & & \multicolumn{4}{p{0.6\linewidth}}{
        \cellcolor[rgb]{0.9,0.9,0.9}{
            \makecell[{{p{\linewidth}}}]{
                \texttt{\tiny{[GM$\rangle$P2]}}
                \texttt{这个分配方案对我不太有利。让我分析一下：} \\
\\ 
\texttt{如果按你的提议：你得2线+1娃娃=10分，我得2铃铛=2分，总分12分。} \\
\\ 
\texttt{但我们可以做得更好：如果**你得2线（8分），我得1娃娃+2铃铛（2+2=4分）**，总分也是12分，但分配更均衡。} \\
\\ 
\texttt{或者：如果你觉得娃娃对你很重要，我们可以考虑其他组合。你愿意接受"你得2线，我得1娃娃+2铃铛"这个方案吗？} \\
            }
        }
    }
    & \\ \\

    \theutterance \stepcounter{utterance}  
    & & & \multicolumn{4}{p{0.6\linewidth}}{
        \cellcolor[rgb]{0.9,0.9,0.9}{
            \makecell[{{p{\linewidth}}}]{
                \texttt{\tiny{[GM$\langle$P2]}}
                \texttt{我理解你想要更均衡的分配。让我计算一下几个可行方案的总分：} \\
\\ 
\texttt{1. **你提议的方案**：我得2线（8分），你得1娃娃+2铃铛（4分） → 总分12分} \\
\texttt{2. **我最初提议**：我得2线+1娃娃（10分），你得2铃铛（2分） → 总分12分} \\
\texttt{3. **另一个方案**：我得1线+1娃娃（6分），你得1线+2铃铛（3+2=5分） → 总分11分} \\
\\ 
\texttt{看来在物品固定的情况下，最大总分是12分。既然总分相同，我同意你的方案：**我得2线，你得1娃娃+2铃铛**。} \\
\\ 
\texttt{{[}提议: 2 线{]}} \\
            }
        }
    }
     \\ \\

    \theutterance \stepcounter{utterance}  
    & & & \multicolumn{2}{p{0.3\linewidth}}{
        \cellcolor[rgb]{0.9,0.9,0.9}{
            \makecell[{{p{\linewidth}}}]{
                \texttt{\tiny{[GM$|$GM]}}
                \texttt{proposal} \\
            }
        }
    }
    & & \\ \\

    \theutterance \stepcounter{utterance}  
    & & & \multicolumn{2}{p{0.3\linewidth}}{
        \cellcolor[rgb]{0.9,0.9,0.9}{
            \makecell[{{p{\linewidth}}}]{
                \texttt{\tiny{[GM$|$GM]}}
                [0, 0, 2]
            }
        }
    }
    & & \\ \\

    \theutterance \stepcounter{utterance}  
    & & \multicolumn{4}{p{0.6\linewidth}}{
        \cellcolor[rgb]{0.9,0.9,0.9}{
            \makecell[{{p{\linewidth}}}]{
                \texttt{\tiny{[P1$\langle$GM]}}
                \texttt{另一位玩家已提交了一个秘密提案。现在轮到你提交你自己的提案。} \\
            }
        }
    }
    & \\ \\

    \theutterance \stepcounter{utterance}  
    & \multicolumn{4}{p{0.6\linewidth}}{
        \cellcolor[rgb]{0.9,0.9,0.9}{
            \makecell[{{p{\linewidth}}}]{
                \texttt{\tiny{[P1$\rangle$GM]}}
                \texttt{{[}提议: 1 娃娃, 2 铃铛{]}} \\
            }
        }
    }
    & & \\ \\

    \theutterance \stepcounter{utterance}  
    & & & \multicolumn{2}{p{0.3\linewidth}}{
        \cellcolor[rgb]{0.9,0.9,0.9}{
            \makecell[{{p{\linewidth}}}]{
                \texttt{\tiny{[GM$|$GM]}}
                \texttt{proposal} \\
            }
        }
    }
    & & \\ \\

    \theutterance \stepcounter{utterance}  
    & & & \multicolumn{2}{p{0.3\linewidth}}{
        \cellcolor[rgb]{0.9,0.9,0.9}{
            \makecell[{{p{\linewidth}}}]{
                \texttt{\tiny{[GM$|$GM]}}
                [2, 1, 0]
            }
        }
    }
    & & \\ \\

    \theutterance \stepcounter{utterance}  
    & & & \multicolumn{2}{p{0.3\linewidth}}{
        \cellcolor[rgb]{0.9,0.9,0.9}{
            \makecell[{{p{\linewidth}}}]{
                \texttt{\tiny{[GM$|$GM]}}
                [[2, 1, 0], [0, 0, 2]]
            }
        }
    }
    & & \\ \\

\end{supertabular}
}

\end{document}
